\section{Square grids and general circuits}\label{sec:grid}

Let $G_n$ be the $n$-by-$n$ square grid. Summation over edge labels
places the family $(Z_{G_n,n})_n$ in $\VNP$. For the circuit lower
bound, we use the fact that every partition
placing at least a third of the vertices on each side cuts many edges.
A small circuit would express the coefficient tensor as a sum in which
each term has low rank across some such partition. The partition may
vary with the term, as permitted in Lemma~\ref{lem:single}.

\begin{theorem}[Grid circuit lower bound]\label{thm:grid}
Let $G$ be the $r$-by-$r$ square grid, where
$r\ge2$ and $r^2\le(\log q)^{1/3}$. Over every field, the state sum
$Z_{G,q}$ has degree-component border circuit size at least
$q^{\Omega(\sqrt r)}$ for all sufficiently large $q$.
\end{theorem}

To obtain the partitions, we decompose the degree-$d$ component of
the circuit output into products of homogeneous factors of balanced
degrees. This is a gate decomposition of the kind used by Valiant,
Skyum, Berkowitz and Rackoff~\cite{VSBR}. Extracting the coefficients with one variable per slot turns each
product into outer products across balanced sets of grid vertices.

\begin{lemma}[Balanced slot sets]\label{lem:balanced}
Suppose the variables are grouped into a set $V$ of $d\ge2$ slots, and
a commutative arithmetic circuit of size $s$ computes $f$. Its slot
tensor has a decomposition
\[
 A(f)=\sum_{\substack{U\subseteq V\\d/3<|U|\le 2d/3}}T_U ,
 \qquad
 \operatorname{rank}_{U\mid V\setminus U}(T_U)
 \le C(s+1)(d+1)^2
\]
with an absolute constant $C$. In particular there are at most $2^d$
terms. For a homogeneous target of degree $d$ and degree-component
border circuit size at most $s$, the slot tensor belongs to the
Zariski closure of the set of such sums.
\end{lemma}

\begin{proof}
Write $[f]_d$ for the degree-$d$ component of $f$. At every gate,
compute all homogeneous components through degree $d$, using
convolution at products. This gives a homogeneous circuit for $[f]_d$
with at most $C(s+1)(d+1)^2$ gates and inputs. The count includes
splitting each of the $O(s+1)$ affine inputs used by the original
circuit into its constant and linear parts.

We maintain an expression for $[f]_d$ as a sum of products $Q_vP_v$,
where $P_v$ is the polynomial at gate $v$ and $Q_v$ is its multiplier.
Initially the only term is the output, with multiplier one. Process
gates in reverse topological order, adding all multipliers at each gate
before processing it. Retain a term when $\deg P_v\le 2d/3$. At a
sum or difference of higher degree, distribute $Q_v$ to both inputs
with the appropriate signs. At a product $P_v=P_uP_t$, choose an
input of larger degree, say $u$, and pass the multiplier $Q_vP_t$ to
$u$. In a tie choose the first input position. This identity works
also for a square in characteristic two. Throughout, take gate degrees
from the homogeneous construction and omit zero terms.

Each step preserves the polynomial represented by the sum. A sum gate
preserves degree along each input, and the larger input of a product
of degree greater than $2d/3$ has degree greater than $d/3$.
Consequently the retained terms give an identity
\[
 [f]_d=\sum_v Q_vP_v,\qquad d/3<\deg P_v\le 2d/3,\qquad
 \deg Q_v=d-\deg P_v .
\]
The number of terms is at most the number of gates and inputs in the
homogeneous circuit, since each retained gate has a single accumulated
multiplier.

For $U\subseteq V$ let $A_U(P)$ denote the array of coefficients of
monomials in $P$ containing one variable from each slot of $U$ and
no other variables. In any contribution to a slot coefficient of
$Q_vP_v$, the variables contributed by $P_v$ occupy some set $U$ of
size $\deg P_v$. The variables contributed by $Q_v$ occupy its
complement. Collecting contributions for each $U$ gives
\[
 T_U=\sum_{v:\,\deg P_v=|U|}
       A_U(P_v)\otimes A_{V\setminus U}(Q_v),
\]
where the tensor factors are placed in the original slot order. Each
summand has rank at most one across this cut, which gives the asserted
rank bound. Finally, $A(f)=A([f]_d)$, and the coefficient map on the
space of homogeneous degree-$d$ polynomials is linear. It therefore
takes the closure defining border size into the closure of the sums
above.
\end{proof}

\begin{proof}[Proof of Theorem~\ref{thm:grid}]
Put $d=r^2$. Color the vertices according to the two sides of a cut in
Lemma~\ref{lem:balanced}. Each color occupies at least $r^2/3$
vertices. If at least $r/3$ rows contain both colors, each of these
rows supplies a horizontal crossing edge. Otherwise the mixed rows
contain fewer than $r^2/3$ vertices altogether. Each color must then
occupy an entire row, so every column supplies a vertical crossing
edge. In both cases the boundary has size at least $\lceil r/3\rceil$.

Let $s$ be the degree-component border circuit size of $Z_{G,q}$.
The grid has $e=2r(r-1)\le 2d$ edges. Lemma~\ref{lem:single} applied
to the closure in Lemma~\ref{lem:balanced}, with
$b=\lceil r/3\rceil$, yields
\[
 \log\!\bigl(C(s+1)(d+1)^2\bigr)
        \ \ge\ \tfrac12\sqrt{\lceil r/3\rceil}\,\log q.
\]
Since $d\le(\log q)^{1/3}$, the contribution of $C(d+1)^2$ on the
left is $O(\log\log q)$. This proves the bound uniformly in $r$.
\end{proof}

For the family in Theorem~\ref{thm:main}, take $q=n$ and restrict the
$n$-by-$n$ grid to an $r$-by-$r$ subgrid with
$r=\lfloor(\log n)^{1/6}\rfloor$. Pinning the unused edges and vertices
as in Section~\ref{sec:notation} gives the subgrid polynomial by
projection. If the full grid has border circuit size $s$, its degree
$n^2$ and Lemma~\ref{lem:projection} give a border circuit of size
$O((s+1)(n^2+1)^2)$ for the subgrid. Theorem~\ref{thm:grid} requires
this size to be at least $n^{\Omega((\log n)^{1/12})}$. The polynomial
overhead of projection is absorbed in the exponent.

It remains to realize the full family as a Boolean sum over any field
$F$. Encode each edge state by a block of $\lceil\log n\rceil$ Boolean
variables. At a vertex, multiply each table entry by the indicator
that its incident blocks encode the corresponding edge labels, and
sum over the entries. Each indicator is a product of bits and their
complements. The product of the vertex polynomials selects exactly
the monomial for a legal edge labeling and vanishes on every unused
encoding. This product has polynomial degree and circuit size since
there are $n^2$ vertices, each of degree at most four and with at most
$n^4$ table entries. Summing it over the Boolean blocks yields
$Z_{G_n,n}$, proving membership in $\VNP$. In the complex setting,
Christandl et al.\ proved $\mathsf{VNP}$-hardness for square-lattice
contractions already with two states per edge~\cite[Theorem~26]{CLSW}.
